% Generated draft: requires mathematical and source review.
\documentclass[11pt]{article}
\usepackage[margin=25mm]{geometry}
\usepackage{fontspec}
\setmainfont{Latin Modern Roman}
\usepackage{amsmath,amssymb,mathtools,unicode-math}
\setmathfont{Latin Modern Math}
\usepackage{enumitem,needspace,xurl,hyperref}
\setlength{\parindent}{0pt}
\setlength{\parskip}{5pt}
\setlength{\emergencystretch}{3em}
\newcommand{\N}{\mathbb{N}}
\newcommand{\Z}{\mathbb{Z}}
\newcommand{\Q}{\mathbb{Q}}
\newcommand{\R}{\mathbb{R}}
\newcommand{\C}{\mathbb{C}}
\newcommand{\F}{\mathbb{F}}
\newcommand{\A}{\mathbb{A}}
\newcommand{\PP}{\mathbb{P}}
\newcommand{\E}{\mathbb{E}}
\newcommand{\eps}{\varepsilon}
\newcommand{\dd}{\,\mathrm{d}}
\begin{document}
\section*{UnsolvedMath ID: 800156}
\textbf{Source problem code:} catalogue entry
\subsection{Definitions and mathematical statement}
The supplied excerpt does not define \(I_2(T)\) or give the coefficients of \(P_4\). The following formulation is conditional on the conventional fourth-moment interpretation of this notation.

Let \(\zeta(s)\) denote the Riemann zeta function, defined for \(\operatorname{Re}s>1\) by \(\zeta(s)=\sum_{n=1}^{\infty}n^{-s}\) and elsewhere by meromorphic continuation. Under the interpretation
\[
I_2(T)=\int_0^T\left|\zeta\!\left(\tfrac12+it\right)\right|^4\,\mathrm{d}t,\qquad T>0,
\]
let \(P_4\) denote the particular degree-four polynomial referred to in the source, and define
\[
E_2(T)=I_2(T)-T P_4(\log T).
\]
The polynomial is the specified main-term polynomial, not an arbitrary polynomial of degree four; its coefficients are not supplied.

The record attributes to Zavorotnyi (1989) the estimate
\[
E_2(T)=O_{\varepsilon}\!\left(T^{2/3+\varepsilon}\right)
\qquad(T\to\infty)
\]
for every \(\varepsilon>0\). Explicitly, for each \(\varepsilon>0\), there are constants \(C_{\varepsilon}>0\) and \(T_{\varepsilon}>0\) such that
\[
|E_2(T)|\le C_{\varepsilon}T^{2/3+\varepsilon}
\qquad\text{for every real }T\ge T_{\varepsilon}.
\]
\textbf{Problem.} Improve this error-term bound. A power-saving improvement would, for example, establish the existence of a fixed real \(\theta<2/3\) such that, for every \(\varepsilon>0\),
\[
E_2(T)=O_{\varepsilon}\!\left(T^{\theta+\varepsilon}\right)
\qquad(T\to\infty).
\]
The source does not specify a target exponent or restrict the requested improvement to this form. The definition and integration limits of \(I_2\), and the exact main-term polynomial, require confirmation from the cited source.
\subsection{Short English statement}
Under the conventional interpretation of \(I_2(T)\) as the fourth moment of the Riemann zeta function on the critical line, improve the error bound \(O_{\varepsilon}(T^{2/3+\varepsilon})\) attributed in the record to Zavorotnyi (1989). The excerpt omits the definition of \(I_2\) and the coefficients of its degree-four main-term polynomial, so this formulation requires confirmation.
\subsection{Sources}
\begin{itemize}
\item Local UnsolvedMath registry, permanent ID \(800156\), source problem code \texttt{catalogue entry}, title: ``Moments of L-functions --- Moments of the Riemann zeta function --- Problem 5.'' Supplied upstream revision: \texttt{d371ee917ac9647fbb06cd6fc529deec6f3dda82}.
\item Hamieh, Kadiri, Ng, \emph{Moments of L-functions Problem List} (2024), Problem 5, as identified by the supplied record: \url{https://arxiv.org/abs/2405.17800}. The supplied provenance reports access on 2026-10-16 and extraction by \texttt{arxiv\_latex}; neither access nor source contents have been independently verified here.
\item The problem text credits Nathan Ng, Lethbridge. The registry lists Hamieh, Kadiri, Ng as proposers and 2024 as the proposed year.
\item The excerpt attributes the displayed error estimate to Zavorotnyi (1989), but supplies no separate bibliographic reference for that result.
\item The registry labels this problem ``open.'' This is a supplied curation label, not an independently verified assessment of the literature.
\end{itemize}
\end{document}
